\documentclass[conference]{IEEEtran}

\IEEEoverridecommandlockouts

\usepackage{amsmath,amssymb}
\usepackage{graphicx}
\usepackage{url}
\usepackage[hidelinks]{hyperref}
\usepackage{fontspec}
\usepackage{xeCJK}

\setmainfont{Noto Serif}
\setsansfont{Noto Sans}
\setmonofont{Noto Sans Mono}
\setCJKmainfont[AutoFakeSlant=0.2]{Noto Serif CJK SC}
\setCJKsansfont[AutoFakeSlant=0.2]{Noto Sans CJK SC}
\setCJKmonofont{Noto Sans Mono CJK SC}
\XeTeXlinebreaklocale "zh"
\XeTeXlinebreakskip = 0pt plus 1pt

\newcommand{\figwide}[4]{%
\begin{figure*}[!tbp]
\centering
\includegraphics[width=#1]{#2}
\vspace{-0.8em}
\caption{#3 #4}
\vspace{-0.8em}
\end{figure*}
}

\newcommand{\figone}[4]{%
\begin{figure}[!tbp]
\centering
\includegraphics[width=#1]{#2}
\vspace{-0.8em}
\caption{#3 #4}
\vspace{-0.8em}
\end{figure}
}

\newcommand{\figwidelabel}[5]{%
\begin{figure*}[!tbp]
\centering
\includegraphics[width=#1]{#2}
\vspace{-0.8em}
\caption{#3 #4}
\label{#5}
\vspace{-0.8em}
\end{figure*}
}

\newcommand{\figonelabel}[5]{%
\begin{figure}[!tbp]
\centering
\includegraphics[width=#1]{#2}
\vspace{-0.8em}
\caption{#3 #4}
\label{#5}
\vspace{-0.8em}
\end{figure}
}

\newcommand{\inlinetable}[3]{%
\refstepcounter{table}\label{#1}
\vspace{0.4em}
\begin{center}
{\footnotesize\textbf{表~\thetable.} #2\par}
\vspace{0.2em}
\footnotesize
#3
\end{center}
\vspace{0.4em}
}

\begin{document}

\title{运行状态约束的气象-功率融合网络用于区域风电集群预测}

\author{\IEEEauthorblockN{待填写作者}
\IEEEauthorblockA{\textit{待填写单位}\\
待填写城市, 中国\\
email@example.com}
}

\maketitle

\begin{abstract}
在区域风电集群 24 h 预测中，未来气象预报是目标窗口内风电变化的主要驱动信息，但预测起点前的真实功率状态同样决定了气象到功率映射的初始条件。仅依赖未来气象的预测模型容易忽略已经观测到的当前出力水平、短期惯性和多场站同步变化，尤其在目标窗口前段产生状态失配误差。为此，本文提出一种运行状态约束的气象-功率融合预测方法，在保持未来气象主驱动地位的同时，将目标窗口前可用的场站级功率状态作为运行状态约束，条件化模型对未来气象记忆的读取。模型首先将后报气象与未来预报气象编码为可查询的时空气象记忆；随后利用多尺度图编码器编码目标窗口前可用的场站级功率序列，形成表征当前场群运行状态的动态功率上下文；最后将功率上下文作为 memory tokens 注入气象记忆，并通过 future-query decoder 在状态条件化后的融合记忆上并行生成 24 h 区域聚合功率预测序列。基于吉林区域 58 个风电场站和 Pangu-Weather 网格化气象预报的实验结果表明，本文方法相较最优对比基线将整体 MAE 降低约 4.1\%，并在 RMSE 和 \(R^2\) 上取得一致改进，同时优于多种网格气象编码基线。逐预测步结果进一步显示，功率状态条件化能够在 24 h 预测窗口内持续增强气象驱动预测的状态适配能力，其中近端预测的改进尤为显著：H1 MAE 相较最优对比基线降低约 42.4\%。这说明目标窗口前可用的真实功率状态不仅提供额外历史信息，更能够作为运行状态约束改善模型对未来气象的读取和解释。
\end{abstract}

\begin{IEEEkeywords}
区域风电集群预测, 风电功率预测, 气象-功率融合, 运行状态约束, 多步预测
\end{IEEEkeywords}

\section{引言}

区域风电集群预测直接服务于电力系统调度、备用容量配置和新能源消纳。与单场站预测相比，区域集群预测关注多个风电场站的聚合出力，既受区域天气过程驱动，也受预测起点前场群真实运行状态、场站间同步变化和短期功率惯性影响。随着网格化数值天气预报和深度气象预报模型的发展，未来气象预报已经成为日内多步风电预测中的关键外生输入。

对于调度场景而言，区域尺度预测更接近实际功率平衡需求。单个场站的局部误差在聚合后可能相互抵消，但区域尺度的天气传播、同步爬坡和限电恢复等现象仍会造成显著预测偏差。因此，区域集群预测不仅需要描述未来风场变化，还需要在预测起点处对当前场群运行状态进行校准。

本文关注业务化区域风电集群 24 h 预测。该任务具有明确的信息可用性约束：未来网格化气象预报覆盖整个目标窗口，是预测未来功率变化的主要驱动；而场站级实测功率只在目标窗口前可用，主要表征当前出力水平、短期惯性和多场站耦合状态。仅使用未来气象的 Weather-only 方法能够捕捉主要天气驱动，但会忽略预测起点前的真实运行状态，尤其容易在预测窗口前段产生状态失配误差。

这一状态失配并不能简单通过增加气象编码器复杂度解决。即使气象预报能够较好描述未来风场，气象到功率的映射仍受当前功率水平、场站运行状态和短时惯性的影响。如果模型只从未来气象场直接回归区域功率，则预测窗口开始阶段容易偏离真实起点；如果将历史功率与未来气象对称拼接，又可能弱化二者在业务时间轴上的不同角色。

已有风电预测研究从风速时序深度模型 \cite{ref1}、网格化空间表示 \cite{ref2}、场群时空关联建模 \cite{ref3,ref4,ref5,ref6,ref7,ref8,ref9,ref10,ref11,ref12}、NWP 校正和多模态融合 \cite{ref13,ref14,ref15,ref16,ref17,ref18} 等角度提升预测性能。然而，许多融合框架仍将历史功率和未来气象作为相对对称的输入，通过拼接、并行编码或特征级注意力进行交互。本文认为，在业务化 24 h 区域风电预测中，更自然的建模方式是非对称气象-功率融合：未来气象作为主驱动记忆，近期功率作为预测起点状态条件，对未来气象信息的读取过程进行约束。

基于这一思想，本文提出一种运行状态约束的气象-功率融合预测网络。模型首先将后报气象和未来气象预报编码为可查询的气象记忆；随后通过多尺度图编码器从最近场站级功率中提取当前场群状态，并生成动态功率上下文 token；最后将这些 token 前置注入气象记忆，使每个 future query 在 cross-attention 中同时读取气象主驱动和功率状态条件，输出 24 步区域集群预测序列。

本文主要贡献如下：
\begin{itemize}
\item 从业务时间可用性出发，将区域风电集群预测表述为非对称气象-功率融合问题，明确未来气象的主驱动作用和近期功率的运行状态约束作用。
\item 提出功率状态条件化的气象-功率融合网络，在 memory 层面注入动态功率上下文 token，而不是仅在输出层进行特征拼接。
\item 在吉林区域 58 个风电场站数据上验证方法有效性。结果表明，本文方法整体优于 Weather-only 和多种网格气象编码基线；逐预测步分析和消融实验说明，真实功率状态能够增强 24 h 预测窗口内的状态适配能力，且起始时段收益尤为显著。
\end{itemize}

\figwide{0.78\textwidth}{../../paper_image_use/method/Fig01_svg.pdf}{业务时间协议。}{近期场站功率仅在目标窗口前可用，未来气象预报覆盖完整预测窗口；本文在该时间可用性约束下进行气象-功率融合预测。}

\figwidelabel{0.78\textwidth}{../../paper_image_use/method/Fig02_svg.pdf}{整体框架。}{未来气象构成主记忆，近期场站功率生成动态上下文 token，并在 memory 层面对气象读取过程进行状态条件化。}{fig:framework}

\section{相关工作与差异}

区域或场群风电预测需要同时刻画时间变化和场站间空间相关性。早期深度学习方法利用 CNN-LSTM 建模风速时序特征 \cite{ref1}，也有工作将风机嵌入规则网格，使卷积网络能够处理空间分布信息 \cite{ref2}。围绕场群时空关联的研究进一步将风电场站间的地理距离、气象相关性、功率相关性和趋势相似性纳入建模，已在层级图建模、动态图建模、趋势感知、空间信息集成和持续学习等场群预测任务中取得进展 \cite{ref3,ref4,ref5,ref6,ref7,ref8,ref9,ref10,ref11,ref12}。这些研究说明，多场站功率状态能够提供有价值的空间耦合信息，但其重点通常是历史功率动态或场群空间结构建模，而不是在未来气象主驱动下进行状态条件化融合。

未来气象预报是多步风电预测中的主要外生信息。已有方法将 NWP 与历史功率进行多模态融合，或围绕 NWP 预报误差、时间滞后、多源时空特征和空间相关性进行建模 \cite{ref13,ref14,ref15,ref16,ref17,ref18}。对于网格化气象场，也有研究使用 3D CNN、区域深度学习框架或概率预测框架直接编码气象数据立方体 \cite{ref19,ref20,ref21}。本文采用 Pangu-Weather 作为未来气象来源 \cite{ref22}，并将其视为预测窗口内的主驱动信息。

Transformer 及其变体为多步时间序列预测提供了并行建模基础。Informer、PatchTST、Crossformer 和 Non-stationary Transformers 分别从长序列注意力、时间片段表示、跨维依赖和非平稳性角度改进序列预测 \cite{ref23,ref24,ref25,ref26}。这些方法说明，注意力机制适合处理多步预测中的长时依赖，但并不直接回答业务化风电预测中未来气象与预测起点运行状态如何分工的问题。

上述两类研究分别强调场站功率时空关系和未来气象驱动，但在业务预测协议下二者并不对称。未来气象在整个预测窗口内可用，而真实功率只在预测窗口前可用。因此，本文并不追求让两类输入以相同方式进入模型，而是让功率状态在气象读取阶段起条件化作用。这一点与常见的多模态分支融合有所区别。

与现有特征拼接、并行分支或对称注意力融合不同，本文关注 memory-level 的非对称融合。query-based decoding 和可查询记忆结构已在结构化输出和通用感知模型中得到应用 \cite{ref27,ref28,ref29}。本文借鉴这一形式，但不进行语言 prompt learning 或外部检索增强；dynamic power context tokens 完全由当前样本的场站级功率动态生成，并用于条件化 future queries 对气象记忆的读取过程。

\section{问题定义与业务时间对齐}

设预测日为 \(D\)，目标窗口为北京时间 \(D\) 日 00:00--23:00。本文直接预测区域级聚合出力序列 \(\mathbf{y}\in\mathbb{R}^{24}\)，即目标区域内 58 个风电场站在每个目标时刻的合计出力，不输出单场站功率。目标窗口前最近 16 步场站级功率作为状态上下文，记为 \(\mathbf{P}^{c}\in\mathbb{R}^{N\times16}\)，其中 \(N=58\) 为场站数量；模型输出为 \(\hat{\mathbf{y}}\in\mathbb{R}^{24}\)。

本文遵循业务化时间可用性约束。对于日样本评价，最新可用功率上下文覆盖 \(D-1\) 日 08:00--23:00。Pangu-Weather 预报在 \(D-1\) 日 08:00 起报，lead 0--15 对应目标窗口之前的后报气象段，lead 16--39 对应 \(D\) 日 00:00--23:00 的未来预报气象段。因此，气象输入包括后报气象 \(\mathbf{X}^{h}\in\mathbb{R}^{16\times C\times H\times W}\) 和未来气象 \(\mathbf{X}^{f}\in\mathbb{R}^{24\times C\times H\times W}\)，二者沿时间维拼接为 \(\mathbf{X}\in\mathbb{R}^{40\times C\times H\times W}\)。场站邻接矩阵 \(\mathbf{A}\in\mathbb{R}^{N\times N}\) 由近期场站功率相关性构造，用于功率状态编码。

训练阶段采用滑动窗口样本，以提高样本数量并覆盖更多预测起点；验证和测试阶段采用严格日样本，以模拟每天固定起报后的业务化预测流程。无论训练还是测试，模型均只使用目标窗口开始前真实可用的信息，避免未来功率泄漏。该设置使本文问题不同于离线随机切分的序列预测任务，也使功率状态条件化的作用能够在逐 horizon 分析中被直接检验。

\section{运行状态约束的气象-功率融合网络}

\subsection{整体框架}

本文方法采用运行状态约束的非对称气象-功率融合设计，如图~\ref{fig:framework} 所示。Weather Encoder 将后报气象和未来预报气象编码为 weather memory tokens，作为预测窗口内的主驱动信息；Power-State Encoder 从目标窗口前可用的场站级功率中提取当前场群状态，并生成 dynamic power context tokens；Query Decoder 使用 24 个可学习 future queries 从状态条件化后的融合记忆中并行读取信息，得到 24 步区域集群预测序列。

该框架的核心不是简单增加功率输入分支，而是改变融合发生的位置。传统拼接式融合通常在编码后或预测头前合并两类特征，模型需要自行学习二者的角色差异。本文将气象特征组织为可查询记忆，并把功率状态转换为同一隐空间中的上下文 token，使每个未来预测步在 cross-attention 读取气象信息时都能显式访问预测起点状态。

因此，功率状态在本文中不是独立的历史预测分支，而是面向未来气象读取过程的运行状态约束。

\subsection{气象记忆构建}

气象记忆构建过程如图~\ref{fig:weather-memory} 所示。Weather Encoder 的输入为 \(\mathbf{X}=[\mathbf{X}^{h};\mathbf{X}^{f}]\)，其中 \(\mathbf{X}^{h}\) 为目标窗口前的后报气象，\(\mathbf{X}^{f}\) 为目标窗口内的未来预报气象。为避免将所有气象变量视为同质通道，本文将输入变量划分为风场变量、背景气象变量和一阶差分变量，并分别通过 frame-wise CNN 提取逐时刻空间特征：
\begin{equation}
\mathbf{F}^{u}_{t}=\phi_{u}(\mathbf{X}^{u}_{t}),\quad
\mathbf{F}^{b}_{t}=\phi_{b}(\mathbf{X}^{b}_{t}),\quad
\mathbf{F}^{\Delta}_{t}=\phi_{\Delta}(\Delta\mathbf{X}_{t}),
\end{equation}
其中 \(u\)、\(b\) 和 \(\Delta\) 分别表示风场、背景场和气象变化分支，\(\phi_{u}\)、\(\phi_{b}\) 和 \(\phi_{\Delta}\) 分别为三条分支的逐帧二维卷积 stem，用于在每个气象时刻提取局部空间特征。风场变量直接影响风电出力，是气象侧的主要可预测性来源；背景场变量刻画大尺度天气形势；一阶差分分支显式描述气象场随 lead time 的变化。

三路特征经 \(1\times1\) 融合后加入时间位置嵌入和阶段嵌入：
\begin{equation}
\mathbf{H}_{t}=\psi([\mathbf{F}^{u}_{t};\mathbf{F}^{b}_{t};\mathbf{F}^{\Delta}_{t}])
+\mathbf{e}^{time}_{t}+\mathbf{e}^{stage}_{t}.
\end{equation}
其中 \(\psi(\cdot)\) 表示 \(1\times1\) 卷积融合块，用于混合三类气象分支的通道信息；\(\mathbf{H}_{t}\) 是第 \(t\) 个气象帧的融合隐表示；\(\mathbf{e}^{time}_{t}\) 为时间位置嵌入，标识气象帧在 40 步输入序列中的相对位置；\(\mathbf{e}^{stage}_{t}\) 为阶段嵌入，用于区分后报气象段和未来预报气象段，避免模型将目标窗口前的气象上下文和目标窗口内的未来预报视为完全同质序列。

随后，将 \(\mathbf{H}_{1:T}\) 输入时空卷积主干 \(\mathcal{E}_{w}\)。该主干由 depthwise 3D 时空卷积、空间下采样卷积和逐帧空间混合模块组成，用于在保留时间长度的同时建模局部时空变化。编码后的特征先经过 \(1\times1\times1\) memory projection，再通过自适应平均池化得到固定大小的时空网格：
\begin{equation}
\mathbf{Z}^{w}=
\operatorname{Pool}\left(\operatorname{Proj}(\mathcal{E}_{w}(\mathbf{H}_{1:T}))\right)
\in\mathbb{R}^{T_m\times H_m\times W_m\times d}.
\end{equation}
其中 \(\operatorname{Proj}\) 为 \(1\times1\times1\) 卷积投影，\(\operatorname{Pool}\) 为自适应平均池化；本文实现中 \(T_m=40\)、\(H_m=6\)、\(W_m=10\)。为保留 memory token 的时空位置，本文将时间、行和列三类可学习位置嵌入相加：
\begin{equation}
\mathbf{E}^{pos}_{\tau,i,j}=
\mathbf{p}^{time}_{\tau}+\mathbf{p}^{row}_{i}+\mathbf{p}^{col}_{j}.
\end{equation}
最终气象记忆定义为
\begin{equation}
\mathbf{M}^{w}=\operatorname{Flatten}(\mathbf{Z}^{w}+\mathbf{E}^{pos})
\in\mathbb{R}^{L_w\times d},\quad L_w=T_mH_mW_m.
\end{equation}
该记忆同时保留目标窗口前的气象上下文、未来气象驱动和局部气象变化信息。

与直接将未来气象压缩为一个全局向量不同，weather memory 保留了时间和空间网格结构，使不同 future queries 可以根据预测 horizon 读取不同 lead time 和空间区域的信息。这一点对于区域集群预测尤为重要，因为多个场站分布在不同位置，区域天气过程在空间上的传播会影响聚合功率变化。

\figonelabel{\columnwidth}{../../paper_image_use/method/Fig03_svg.pdf}{气象记忆构建。}{解耦气象分支提取风场、背景场和气象变化特征，随后构造可被 future queries 读取的时空气象记忆。}{fig:weather-memory}

\subsection{功率状态编码}

功率状态编码过程如图~\ref{fig:power-state} 所示。功率状态编码器输入目标窗口前最近 \(T_c\) 步场站级功率 \(\mathbf{P}^{c}\in\mathbb{R}^{N\times T_c}\) 和场站邻接矩阵 \(\mathbf{A}\in\mathbb{R}^{N\times N}\)，其中 \(N\) 为场站数量，本文实现中 \(T_c=16\)、\(N=58\)。\(\mathbf{P}^{c}\) 是预测起点前真实可用的场站功率状态，不包含目标窗口内的未来功率。图卷积使用自环增强和对称归一化：
\begin{equation}
\tilde{\mathbf{A}}=\mathbf{A}+\mathbf{I},\quad
\hat{\mathbf{A}}=\tilde{\mathbf{D}}^{-1/2}\tilde{\mathbf{A}}\tilde{\mathbf{D}}^{-1/2},
\end{equation}
\begin{equation}
\operatorname{GCN}(\mathbf{H})=\hat{\mathbf{A}}\mathbf{H}\mathbf{W}+\mathbf{b}.
\end{equation}

模型在 point、patch 和 global 三种时间尺度上编码场站功率。point 尺度逐时刻处理 \(\mathbf{P}^{c}_{:,t}\)，用于刻画每个已知时刻的场站空间状态；patch 尺度处理长度为 \(s\)、步长为 \(r\) 的局部功率片段，用于刻画短时爬坡或回落模式；global 尺度处理完整 \(\mathbf{P}^{c}\)，用于刻画整个上下文窗口的出力水平和整体趋势。三类 token 可写为
\begin{equation}
\mathbf{o}^{pt}_{t}=g_{pt}(\operatorname{GCN}_{pt}(\mathbf{P}^{c}_{:,t},\mathbf{A})),
\end{equation}
\begin{equation}
\begin{aligned}
\mathbf{o}^{pa}_{j}
&=g_{pa}(\operatorname{GCN}_{pa}(\mathbf{P}^{c}_{:,jr:jr+s},\mathbf{A})),\\
\mathbf{o}^{gl}
&=g_{gl}(\operatorname{GCN}_{gl}(\mathbf{P}^{c},\mathbf{A})).
\end{aligned}
\end{equation}
其中 \(g_{pt}\)、\(g_{pa}\) 和 \(g_{gl}\) 为展平后的线性投影，分别将图卷积输出映射到 \(d\) 维 token 空间；\(s\) 和 \(r\) 分别表示局部片段长度和滑动步长，patch token 数量为 \(J=\lfloor (T_c-s)/r\rfloor+1\)。

为使三种尺度的信息相互校准，功率编码器采用自底向上和自顶向下的跨尺度注意力。首先，patch tokens 读取 point tokens，聚合逐时刻细粒度变化；随后 global token 读取 patch tokens，形成窗口级运行状态。接着，global token 反向调制 patch tokens，patch tokens 再调制 point tokens，使局部片段和逐时刻表示都带有整体运行趋势。跨尺度融合后的 token 序列记为
\begin{equation}
\mathbf{S}^{p}=[\mathbf{o}^{gl};\mathbf{o}^{pa}_{1:J};\mathbf{o}^{pt}_{1:T_c}]\in\mathbb{R}^{(1+J+T_c)\times d}.
\end{equation}
该序列加入位置编码后输入轻量 Transformer encoder，并经投影得到样本级功率状态向量：
\begin{equation}
\mathbf{z}^{p}=f_{p}(\operatorname{Transformer}(\mathbf{S}^{p}))
\in\mathbb{R}^{d}.
\end{equation}
其中 \(J\) 为 patch 数量，\(f_p\) 为投影头。由于预测目标是区域聚合出力，功率编码器不直接输出场站级预测，而是将当前出力水平、局部变化和场站间耦合关系压缩为样本特定的运行状态表示，供后续气象读取过程使用。

\figonelabel{\columnwidth}{../../paper_image_use/method/Fig04_svg.pdf}{功率状态编码。}{多尺度 GCN 从近期场站功率中提取当前运行状态、局部变化和多场站耦合关系。}{fig:power-state}

\subsection{Memory-Level 融合与查询解码}

本文的关键融合机制是将功率状态写入气象记忆，而不是在输出层前简单拼接特征。Power-to-context adapter 将功率状态向量映射为 \(K\) 个动态功率上下文 token：
\begin{equation}
\mathbf{C}^{p}=\operatorname{reshape}(a_{\theta}(\mathbf{z}^{p}))\in\mathbb{R}^{K\times d}.
\end{equation}
随后将其前置到气象记忆：
\begin{equation}
\mathbf{M}^{pc}=[\mathbf{C}^{p};\mathbf{M}^{w}]\in\mathbb{R}^{(K+L_w)\times d}.
\end{equation}

这种前置注入使功率状态以 memory token 的形式参与所有预测步的 cross-attention。不同 horizon 的 query 可以在同一融合记忆中自适应分配对功率上下文和未来气象 memory 的注意力：起始预测步可利用功率上下文校准初始运行状态，中后段预测步仍保留对未来气象演变的主驱动读取。因此，该结构不是把功率状态限定为单步校正项，而是将其作为贯穿 24 h 解码过程的运行状态约束。

Query Decoder 维护 \(24\) 个 learnable future queries。第 \(\ell\) 层 decoder block 通过 query self-attention 建模不同预测步之间的依赖，并通过 cross-attention 读取融合记忆：
\begin{equation}
\mathbf{Q}^{(\ell)}=\operatorname{DecoderBlock}(\mathbf{Q}^{(\ell-1)},\mathbf{M}^{pc}).
\end{equation}
最终预测为：
\begin{equation}
\hat{\mathbf{y}}=g_{\phi}(\operatorname{LN}(\mathbf{Q}^{(L_d)})).
\end{equation}
训练仅使用最终 24 步区域聚合功率预测损失：
\begin{equation}
\mathcal{L}=\lambda_1\,\operatorname{MAE}(\hat{\mathbf{y}},\mathbf{y})
+\lambda_2\,\operatorname{RMSE}(\hat{\mathbf{y}},\mathbf{y}).
\end{equation}
\(\lambda_1\) 和 \(\lambda_2\) 用于平衡绝对误差和平方误差项，本文取 \(\lambda_1=0.6\)、\(\lambda_2=0.4\)。

所有 future queries 并行解码，避免自回归误差累积；训练目标直接对应最终 24 h 区域风电预测任务，便于在相同监督信号下比较不同融合结构。

\section{实验}

\subsection{数据与设置}

本文基于吉林省区域风电场群构建实验数据集，预测对象为 58 个风电场站形成的区域聚合出力，并配套使用覆盖研究区域的 Pangu-Weather 网格化气象预报。气象网格大小为 \(19\times33\)，输入通道数为 17，包括 10 m 风场、不同气压层风场、温度、湿度、位势高度和海平面气压等变量。所有功率和气象数据均按北京时间对齐。

训练集使用 2022 年滑动窗口样本，以增加训练样本数量并覆盖更多天气和运行状态；验证集和测试集使用 2023 年严格日样本，其中每月前 10 天作为验证集，其余日期作为测试集。所有模型采用相同的数据划分、归一化方式和 24 h 预测目标。评价指标包括归一化 MAE、RMSE 和 \(R^2\)。

本文方法使用目标窗口前 16 步场站级功率作为运行状态上下文。

对比模型包括 Weather-only 以及多种网格气象编码基线。Weather-only 与本文方法共享气象输入和 query decoder，但不使用功率状态条件，用于直接检验功率状态融合的增益。CNN、3D CNN、ViT、Swin 和 Video Transformer 均以同一气象输入为基础，用于比较不同网格气象编码方式对区域风电预测的影响。

为保证时间信息使用一致，所有 weather-only 基线均输入完整的 16 步后报气象和 24 步未来预报气象。CNN、ViT 和 Swin 基线采用逐帧空间编码后接 temporal Transformer 的方式建模 40 帧时间依赖；3D CNN 和 Video Transformer 则直接进行时空编码。所有基线最终均从未来 24 个时间位置输出区域聚合功率预测。

\subsection{整体预测性能}

表~\ref{tab:benchmark-short} 给出了不同模型的整体测试结果。本文方法在 MAE、RMSE 和 \(R^2\) 上均取得最优表现，分别为 0.0767、0.1041 和 0.8330；对应实际功率尺度下的 MAE(MW) 和 RMSE(MW) 分别为 377.9 MW 和 512.7 MW。Weather-only 是所有基线中最强的模型，说明未来气象预报已经提供了主要预测信号。本文方法在 Weather-only 基础上进一步降低误差，说明近期功率状态并非替代未来气象，而是为气象驱动预测提供预测起点状态约束。

\inlinetable{tab:benchmark-short}{整体预测性能}{
\begin{tabular}{lccc}
\hline
Model & MAE \(\downarrow\) & RMSE \(\downarrow\) & \(R^2\) \(\uparrow\) \\
\hline
Ours & \textbf{0.0767} & \textbf{0.1041} & \textbf{0.8330} \\
Weather-only & 0.0800 & 0.1072 & 0.8230 \\
3D CNN & 0.0845 & 0.1138 & 0.8003 \\
ViT & 0.0899 & 0.1193 & 0.7806 \\
Video Trans. & 0.0909 & 0.1211 & 0.7738 \\
CNN & 0.0926 & 0.1226 & 0.7684 \\
Swin & 0.0947 & 0.1255 & 0.7571 \\
\hline
\end{tabular}
}

除 Weather-only 外，3D CNN 是误差最低的网格气象编码基线，但仍明显弱于本文方法。ViT、Swin 和 Video Transformer 在该数据集上未超过 Weather-only，说明直接替换更复杂的气象视觉编码器并不能自动解决预测起点状态失配问题。相较之下，本文的提升来自对功率状态和未来气象角色的重新划分：气象负责提供预测窗口内主驱动，功率状态负责修正起点附近运行状态。

从指标差异看，本文方法相对 Weather-only 的整体 MAE 降幅约为 4.1\%，RMSE 降幅约为 2.9\%。由于 Weather-only 已经使用完整未来气象预报，该提升并非来自额外未来信息，而来自目标窗口前可用的真实功率状态。该结果也解释了为什么本文不把功率状态建模为独立预测器，而是将其作为条件信息注入气象读取过程。

\subsection{逐预测步融合效果}

表~\ref{tab:horizon-short} 和图~\ref{fig:horizon} 展示了本文方法与 Weather-only 在不同预测步长上的误差差异，其中 H\(n\) 对应第 \(n\) 个预测步。H1 上本文方法的 MAE 为 0.0430，明显低于 Weather-only 的 0.0747；RMSE 也从 0.0956 降至 0.0586。这说明预测起点前的场站功率状态对目标窗口开始阶段具有强约束作用，能够显著缓解 Weather-only 在近端的状态失配。

\inlinetable{tab:horizon-short}{逐预测步 MAE 对比}{
\begin{tabular}{lccccc}
\hline
Model & H1 & H6 & H12 & H24 & Avg \\
\hline
Ours & \textbf{0.0430} & \textbf{0.0722} & \textbf{0.0786} & 0.0903 & \textbf{0.0767} \\
Weather & 0.0747 & 0.0748 & 0.0810 & \textbf{0.0886} & 0.0800 \\
\hline
\end{tabular}
}

\figonelabel{\columnwidth}{../../paper_image_use/result/fig09_ours_weather_horizon_residual.pdf}{逐预测步误差及差值。}{本文模型降低 24 步平均误差，起始预测步相对 Weather-only 的改善尤为明显。}{fig:horizon}

逐步长曲线显示，状态条件化带来的收益在目标窗口起始阶段最为突出，并进一步带动 24 步平均 MAE 从 0.0800 降至 0.0767。H6 和 H12 上本文方法仍保持优势，H24 上两者误差接近，说明模型没有用功率历史替代未来气象，而是在气象主驱动基础上引入预测起点的运行状态约束。对于调度应用而言，这种起始状态适配能力能够提高日内预测序列与真实运行过程的衔接性。

\subsection{消融实验}

各消融变体用于分别检验功率状态、功率编码结构和气象记忆的作用。Weather-only 移除功率状态上下文，仅使用气象输入；Random power 和 Zero power 保留功率侧结构，但分别用随机功率和零功率替代真实场站功率；Basic GCN 将多尺度功率编码器替换为基础图卷积结构。No memory 关闭时空气象 memory，No wind 移除风场相关气象变量。

表~\ref{tab:ablation-short} 给出了关键消融结果。Weather-only 移除所有功率状态条件，MAE 增加 4.27\%。Random power 和 Zero power 保留功率侧结构但移除真实状态信息，其性能接近 Weather-only，说明完整模型的收益来自真实功率状态，而不是额外分支或参数量。Basic GCN 弱于完整模型，说明多尺度图编码和跨尺度交互有助于形成更有效的场群状态表示。

\inlinetable{tab:ablation-short}{关键消融结果}{
\begin{tabular}{lccc}
\hline
Variant & MAE \(\downarrow\) & RMSE \(\downarrow\) & \(\Delta\)MAE \\
\hline
Ours & \textbf{0.0767} & \textbf{0.1041} & 0.00\% \\
Weather-only & 0.0800 & 0.1072 & 4.27\% \\
Random power & 0.0804 & 0.1079 & 4.72\% \\
Zero power & 0.0804 & 0.1078 & 4.70\% \\
Basic GCN & 0.0785 & 0.1062 & 2.29\% \\
No memory & 0.0808 & 0.1093 & 5.33\% \\
No wind & 0.1518 & 0.2047 & 97.79\% \\
\hline
\end{tabular}
}

气象侧消融进一步说明，本文方法仍以未来气象为主驱动。No memory 关闭气象空间记忆后误差明显上升，说明时空气象 memory 对 query 解码有效；No wind 屏蔽风场相关变量后误差接近翻倍，表明风场变量是该任务中不可替代的主要气象信息。上述结果共同支持本文的非对称融合设定：气象提供未来主驱动，功率提供起点状态条件。

结合逐预测步结果可以看到，功率侧和气象侧消融回答的是不同问题。随机或零功率退化说明模型确实依赖真实运行状态；No memory 和 No wind 退化则说明气象记忆仍是主信息载体。二者共同限定了本文方法的定位：它不是纯历史功率模型，也不是纯气象模型，而是面向业务时间约束的状态条件化融合预测模型。

\section{讨论与结论}

实验结果表明，在区域风电集群 24 h 预测中，未来气象预报仍是主导信息来源。Weather-only 已经强于多种网格气象编码基线，说明 Pangu-Weather 网格化预报能够捕捉预测窗口内的主要外生驱动。本文方法的作用不是用历史功率替代未来气象，而是在气象驱动预测中引入预测起点前的真实场群状态约束，使模型能够在读取未来气象时同时感知当前运行状态。

逐预测步结果进一步说明，运行状态约束的收益具有明确的业务含义：它在完整 24 h 平均指标上带来稳定改善，并在目标窗口起始时段表现最为突出。H1 MAE 的大幅降低表明，当前出力水平、短期惯性和多场站同步变化能够有效修正气象到功率映射的初始偏移；同时，24 步平均 MAE 和 RMSE 的同步下降说明这种约束会通过 query 解码过程影响完整预测序列，而不是只作用于单个时刻。

消融实验进一步区分了两类信息来源。No wind 和 No memory 的明显退化表明，风场变量和时空气象记忆是模型的主信息载体；Random power 和 Zero power 的退化则表明，功率侧收益来自真实运行状态，而不是额外分支或参数量。由此可见，本文方法的关键在于非对称融合：未来气象提供目标窗口内的主驱动，近期场站功率提供预测起点的运行状态约束，二者在 memory-level fusion 中共同服务于区域风电预测。

本文提出了运行状态约束的气象-功率融合网络，将未来气象编码为主记忆，将近期场站级功率编码为动态上下文 token，并通过 memory-level fusion 条件化 future queries 的读取过程。吉林区域 58 个风电场站案例验证了该方法在整体预测误差和起始时段预测误差上的有效性。未来工作可进一步结合气象预报不确定性、跨区域迁移和极端爬坡事件建模，提升模型在复杂运行场景下的适应性。

\begin{thebibliography}{29}

\bibitem{ref1}
H. Liu, X. Mi, and Y. Li, ``Smart deep learning based wind speed prediction model using wavelet packet decomposition, convolutional neural network and convolutional long short term memory network,'' \textit{Energy Conversion and Management}, vol. 166, pp. 120--131, Jun. 2018, doi: \url{https://doi.org/10.1016/j.enconman.2018.04.021}.

\bibitem{ref2}
R. Yu, Z. Liu, X. Li, W. Lu, D. Ma, M. Yu, J. Wang, and B. Li, ``Scene learning: Deep convolutional networks for wind power prediction by embedding turbines into grid space,'' \textit{Applied Energy}, vol. 238, pp. 249--257, Mar. 2019, doi: \url{https://doi.org/10.1016/j.apenergy.2019.01.010}.

\bibitem{ref3}
L. Wang, J. Gao, Y. Li, and D. Wang, ``A method for ultra-short-term wind power forecasting of large-scale wind farms based on adaptive spatiotemporal graph convolution,'' \textit{Renewable Energy}, vol. 249, Art. no. 123321, Aug. 2025, doi: \url{https://doi.org/10.1016/j.renene.2025.123321}.

\bibitem{ref4}
M. Li, M. Yang, and Y. Yu, ``Short-term regional wind power forecast method based on hierarchical graph neural network and extreme value theory,'' \textit{Energy Conversion and Management}, vol. 341, Art. no. 120002, Oct. 2025, doi: \url{https://doi.org/10.1016/j.enconman.2025.120002}.

\bibitem{ref5}
X. Liu, Y. Zhang, Z. Zhen, F. Xu, F. Wang, and Z. Mi, ``Spatio-Temporal Graph Neural Network and Pattern Prediction Based Ultra-Short-Term Power Forecasting of Wind Farm Cluster,'' \textit{IEEE Transactions on Industry Applications}, vol. 60, no. 1, pp. 1794--1803, Jan. 2024, doi: \url{https://doi.org/10.1109/TIA.2023.3321267}.

\bibitem{ref6}
M. Yang, C. Ju, Y. Huang, Y. Guo, and M. Jia, ``Short-Term Power Forecasting of Wind Farm Cluster Based on Global Information Adaptive Perceptual Graph Convolution Network,'' \textit{IEEE Transactions on Sustainable Energy}, vol. 15, no. 3, pp. 2063--2076, Jul. 2024, doi: \url{https://doi.org/10.1109/TSTE.2024.3397877}.

\bibitem{ref7}
M. Yang, D. Wang, W. Zhang, and X. Yv, ``A centralized power prediction method for large-scale wind power clusters based on dynamic graph neural network,'' \textit{Energy}, vol. 310, Art. no. 133210, Nov. 2024, doi: \url{https://doi.org/10.1016/j.energy.2024.133210}.

\bibitem{ref8}
M. Yang, Y. Guo, T. Huang, F. Fan, C. Ma, and G. Fang, ``Wind farm cluster power prediction based on graph deviation attention network with learnable graph structure and dynamic error correction during load peak and valley periods,'' \textit{Energy}, vol. 312, Art. no. 133645, Dec. 2024, doi: \url{https://doi.org/10.1016/j.energy.2024.133645}.

\bibitem{ref9}
M. Yang, Y. Huang, Y. Guo, W. Zhang, and B. Wang, ``Ultra-short-term wind farm cluster power prediction based on FC-GCN and trend-aware switching mechanism,'' \textit{Energy}, vol. 290, Art. no. 130238, Mar. 2024, doi: \url{https://doi.org/10.1016/j.energy.2024.130238}.

\bibitem{ref10}
Z. Qu, J. Li, X. Hou, and J. Gui, ``A D-stacking dual-fusion, spatio-temporal graph deep neural network based on a multi-integrated overlay for short-term wind-farm cluster power multi-step prediction,'' \textit{Energy}, vol. 281, Art. no. 128289, Oct. 2023, doi: \url{https://doi.org/10.1016/j.energy.2023.128289}.

\bibitem{ref11}
M. Yang, Y. Jiang, C. Xu, B. Wang, Z. Wang, and X. Su, ``Day-ahead wind farm cluster power prediction based on trend categorization and spatial information integration model,'' \textit{Applied Energy}, vol. 388, Art. no. 125580, Jun. 2025, doi: \url{https://doi.org/10.1016/j.apenergy.2025.125580}.

\bibitem{ref12}
Q. Zhang, K. Mei, Q. Zhou, X. Ding, Z. Dong, M. Tian, D. Liu, W. He, Z. Zeng, and M. Li, ``A wind farm cluster power prediction method based on globally aware dynamic spatiotemporal graph continual learning framework,'' \textit{Energy}, vol. 348, Art. no. 140591, Apr. 2026, doi: \url{https://doi.org/10.1016/j.energy.2026.140591}.

\bibitem{ref13}
L. Wang, Y. He, X. Liu, L. Li, and K. Shao, ``M2TNet: Multi-modal multi-task Transformer network for ultra-short-term wind power multi-step forecasting,'' \textit{Energy Reports}, vol. 8, pp. 7628--7642, Nov. 2022, doi: \url{https://doi.org/10.1016/j.egyr.2022.05.290}.

\bibitem{ref14}
L. Wang and Y. He, ``M2STAN: Multi-modal multi-task spatiotemporal attention network for multi-location ultra-short-term wind power multi-step predictions,'' \textit{Applied Energy}, vol. 324, Art. no. 119672, Oct. 2022, doi: \url{https://doi.org/10.1016/j.apenergy.2022.119672}.

\bibitem{ref15}
C. Liu, X. Zhang, S. Mei, Q. Zhou, and H. Fan, ``Series-wise attention network for wind power forecasting considering temporal lag of numerical weather prediction,'' \textit{Applied Energy}, vol. 336, Art. no. 120815, Apr. 2023, doi: \url{https://doi.org/10.1016/j.apenergy.2023.120815}.

\bibitem{ref16}
M. Yang, Y. Guo, T. Huang, and W. Zhang, ``Power prediction considering NWP wind speed error tolerability: A strategy to improve the accuracy of short-term wind power prediction under wind speed offset scenarios,'' \textit{Applied Energy}, vol. 377, Art. no. 124720, Jan. 2025, doi: \url{https://doi.org/10.1016/j.apenergy.2024.124720}.

\bibitem{ref17}
M. Yang, C. Han, W. Zhang, and B. Wang, ``A short-term power prediction method for wind farm cluster based on the fusion of multi-source spatiotemporal feature information,'' \textit{Energy}, vol. 294, Art. no. 130770, May 2024, doi: \url{https://doi.org/10.1016/j.energy.2024.130770}.

\bibitem{ref18}
S. Hu, Y. Xiang, H. Zhang, S. Xie, J. Li, C. Gu, W. Sun, and J. Liu, ``Hybrid forecasting method for wind power integrating spatial correlation and corrected numerical weather prediction,'' \textit{Applied Energy}, vol. 293, Art. no. 116951, Jul. 2021, doi: \url{https://doi.org/10.1016/j.apenergy.2021.116951}.

\bibitem{ref19}
K. Higashiyama, Y. Fujimoto, and Y. Hayashi, ``Feature Extraction of NWP Data for Wind Power Forecasting Using 3D-Convolutional Neural Networks,'' \textit{Energy Procedia}, vol. 155, pp. 350--358, Nov. 2018, doi: \url{https://doi.org/10.1016/j.egypro.2018.11.043}.

\bibitem{ref20}
J. Keisler and E. Le Naour, ``WindDragon: automated deep learning for regional wind power forecasting,'' \textit{Environmental Data Science}, vol. 4, Art. no. e19, 2025, doi: \url{https://doi.org/10.1017/eds.2025.10}.

\bibitem{ref21}
J. Jonkers, D. N. Avendano, G. Van Wallendael, and S. Van Hoecke, ``A novel day-ahead regional and probabilistic wind power forecasting framework using deep CNNs and conformalized regression forests,'' \textit{Applied Energy}, vol. 361, Art. no. 122900, May 2024, doi: \url{https://doi.org/10.1016/j.apenergy.2024.122900}.

\bibitem{ref22}
K. Bi, L. Xie, H. Zhang, C. Chen, X. Gu, and Q. Tian, ``Accurate medium-range global weather forecasting with 3D neural networks,'' \textit{Nature}, vol. 619, no. 7970, pp. 533--538, Jul. 2023, doi: \url{https://doi.org/10.1038/s41586-023-06185-3}.

\bibitem{ref23}
H. Zhou, S. Zhang, J. Peng, S. Zhang, J. Li, H. Xiong, and W. Zhang, ``Informer: Beyond Efficient Transformer for Long Sequence Time-Series Forecasting,'' in \textit{Proc. AAAI Conf. Artificial Intelligence}, vol. 35, no. 12, pp. 11106--11115, 2021, doi: \url{https://doi.org/10.1609/aaai.v35i12.17325}.

\bibitem{ref24}
Y. Nie, N. H. Nguyen, P. Sinthong, and J. Kalagnanam, ``A Time Series is Worth 64 Words: Long-term Forecasting with Transformers,'' in \textit{Proc. International Conference on Learning Representations (ICLR)}, 2023. [Online]. Available: \url{https://openreview.net/forum?id=Jbdc0vTOcol}

\bibitem{ref25}
Y. Zhang and J. Yan, ``Crossformer: Transformer Utilizing Cross-Dimension Dependency for Multivariate Time Series Forecasting,'' in \textit{Proc. International Conference on Learning Representations (ICLR)}, 2023. [Online]. Available: \url{https://openreview.net/forum?id=vSVLM2j9eie}

\bibitem{ref26}
Y. Liu, H. Wu, J. Wang, and M. Long, ``Non-stationary Transformers: Exploring the Stationarity in Time Series Forecasting,'' in \textit{Advances in Neural Information Processing Systems}, vol. 35, pp. 9881--9893, 2022. [Online]. Available: \url{http://papers.nips.cc/paper_files/paper/2022/hash/4054556fcaa934b0bf76da52cf4f92cb-Abstract-Conference.html}

\bibitem{ref27}
N. Carion, F. Massa, G. Synnaeve, N. Usunier, A. Kirillov, and S. Zagoruyko, ``End-to-End Object Detection with Transformers,'' in \textit{Computer Vision -- ECCV 2020}, ser. Lecture Notes in Computer Science, vol. 12346. Cham, Switzerland: Springer, 2020, pp. 213--229, doi: \url{https://doi.org/10.1007/978-3-030-58452-8_13}.

\bibitem{ref28}
A. Jaegle, S. Borgeaud, J.-B. Alayrac, C. Doersch, C. Ionescu, D. Ding, S. Koppula, D. Zoran, A. Brock, E. Shelhamer, O. J. H{\'{e}}naff, M. M. Botvinick, A. Zisserman, O. Vinyals, and J. Carreira, ``Perceiver IO: A General Architecture for Structured Inputs \& Outputs,'' in \textit{Proc. International Conference on Learning Representations (ICLR)}, 2022. [Online]. Available: \url{https://openreview.net/forum?id=fILj7WpI-g}

\bibitem{ref29}
J. Weston, S. Chopra, and A. Bordes, ``Memory Networks,'' in \textit{Proc. International Conference on Learning Representations (ICLR)}, 2015. [Online]. Available: \url{http://arxiv.org/abs/1410.3916}

\end{thebibliography}


\end{document}
